% Dalies P3 išsamus funkcinių taškų skaičiavimas (atsakingas Marijus Planciunas).
% Įterpiama priede po 87 funkcijų lentelės; skaičiai turi sutapti su P3 eilutėmis tab:fp-sarasas lentelėje.

\subsection{Išsamus dalies P3 skaičiavimas}\label{sec:priedas-p3}

Dalis P3 „Dalyvių valdymas“ (13 funkcijų, 66~UFP) aprašyta išsamiai, nes
joje yra visi keturi sistemoje esantys funkcijų tipai~– ILF, EI, EO ir EQ (EIF sistemoje nėra). Kiekvienai
funkcijai \ref{tab:fp-p3}~lentelėje įvardyti jos DET ir RET arba FTR, o
paskutiniame stulpelyje~– matricos intervalai, kurie lėmė sudėtingumą.

Transakcijoje DET~– naudotojui atpažįstamas laukas, kertantis sistemos
ribą; ILF DET~– naudotojui atpažįstamas saugomas laukas, kurį įrašo arba
skaito bent viena funkcija. DET skaičiuojami pagal IFPUG taisykles: įvestyje
vienas DET skiriamas veiksmo mygtukui ir vienas~– sistemos pranešimui
(patvirtinimui arba klaidai), o prisijungusio naudotojo tapatybė DET
nelaikoma. Užklausas ir išvestis paleidžia pasirinkimas (filtras,
laikotarpis, paieškos mygtukas arba sąrašo įrašas), kuris, kaip ir
navigacijos mygtukas paskaitos pavyzdyje, laikomas veiksmo DET;
priminimą (P3.11) sistema siunčia pati nustatytu laiku. Paskaitos parkavimo
pavyzdyje pranešimas atskiru DET nelaikomas; taip skaičiuojant, lankomumo
pažymėjimas (P3.7, 4~DET ir 2~FTR) taptų žemo sudėtingumo, o P3 būtų
65~UFP. Vardas ir pavardė yra du DET, kai rodomi atskiruose stulpeliuose
ar kortelės laukuose.
RET~– loginis ILF pogrupis, FTR~– ILF, kurį transakcija skaito arba keičia.
FTR įvardyti ILF pavadinimais iš \ref{tab:fp-sarasas}~lentelės: „Naudotojo
paskyra“ (P0.1), „Sporto erdvė“ (P0.2), „Užsiėmimas“ (P2.1) ir
„Rezervacijos ir lankomumas“ (P3.1); skliaustuose nurodyta, ką transakcija
tame faile skaito arba keičia.

\begingroup
\small
\setlength{\tabcolsep}{4pt}
\setlength{\LTcapwidth}{\textwidth}
\captionsetup{font=normalsize}
% 5 stulpeliai, @{} kraštuose – 8 tarpai \tabcolsep; DET stulpelis užpildo likusį plotį.
\begin{longtable}{@{}>{\raggedright\arraybackslash}p{0.85cm}>{\raggedright\arraybackslash}p{2.75cm}>{\raggedright\arraybackslash}p{\dimexpr\textwidth-10.25cm-8\tabcolsep\relax}>{\raggedright\arraybackslash}p{4.3cm}>{\raggedright\arraybackslash}p{2.35cm}@{}}
\caption{Dalies P3 funkcijų DET, RET ir FTR}\label{tab:fp-p3}\\
\toprule
\textbf{Nr.} & \textbf{Funkcija ir tipas} & \textbf{DET} & \textbf{RET arba FTR} & \textbf{Sudėtingumas ir UFP} \\
\midrule
\endfirsthead
\caption[]{Dalies P3 funkcijų DET, RET ir FTR (tęsinys)}\\
\toprule
\textbf{Nr.} & \textbf{Funkcija ir tipas} & \textbf{DET} & \textbf{RET arba FTR} & \textbf{Sudėtingumas ir UFP} \\
\midrule
\endhead
\endfoot
\midrule
\multicolumn{5}{@{}l@{}}{\textbf{Iš viso: 13 funkcijų, 66 UFP}} \\
\bottomrule
\endlastfoot
P3.1 & Rezervacijos ir lankomumas\newline ILF
 & \textbf{12}\newline rezervacija (8): užsiėmimas ir dalyvis (kas į kurį
   užsiėmimą užsiregistravo), rezervavimo laikas, būsena (aktyvi, atšaukta,
   pašalinta), kas užregistravo (dalyvis, treneris ar administratorius),
   atšaukimo ar pašalinimo laikas, kas atšaukė ar pašalino, priminimas
   išsiųstas;\newline
   lankomumo įrašas (4): lankomumo žyma (dalyvavo arba nedalyvavo), kas
   pažymėjo, pažymėjimo laikas, pataisymo laikas\par\smallskip
   \footnotesize Laukus įrašo: kas užregistravo~– P3.3, P3.5; atšaukimo ar
   pašalinimo laikas, kas atšaukė ar pašalino~– P3.4, P3.6, P0.12, P2.5;
   priminimas išsiųstas~– P3.11; kas pažymėjo, pažymėjimo ir pataisymo
   laikas~– P3.7.
 & \textbf{RET 2}: rezervacija; lankomumo įrašas (atsiranda tik treneriui
   pažymėjus lankomumą, P3.7)
 & žemas, 7\newline{\footnotesize DET 1–19\newline RET 2–5} \\
\midrule
P3.2 & Užsiėmimų tvarkaraštis su laisvų vietų skaičiumi\newline EO
 & \textbf{10}: laikotarpis (savaitė), užsiėmimo pavadinimas, režimas (su
   treneriu ar virtualus treneris), data, pradžios laikas, pabaigos laikas,
   vieta (pagal ją galima ir filtruoti), treneris, vietų skaičius, laisvų
   vietų skaičius (apskaičiuojamas)
 & \textbf{FTR 4}: Užsiėmimas (pavadinimas, režimas, data ir laikas, vietų
   skaičius); Rezervacijos ir lankomumas (aktyvios rezervacijos); Sporto
   erdvė (vieta); Naudotojo paskyra (treneris)
 & aukštas, 7\newline{\footnotesize DET 6–19\newline FTR 4+} \\
\midrule
P3.3 & Vietos rezervavimas (tikrinant laisvas vietas)\newline EI
 & \textbf{3}: pasirinktas užsiėmimas, mygtukas „Rezervuoti“, pranešimas
   (patvirtinimas arba „Laisvų vietų nebėra“)
 & \textbf{FTR 3}: Užsiėmimas (vietų skaičius, ar dar neprasidėjo);
   Rezervacijos ir lankomumas (užimtos vietos, nauja rezervacija, kas
   užregistravo); Naudotojo
   paskyra (ar paskyra nedeaktyvuota: deaktyvuoto dalyvio rezervacijos
   negalimos, plg. P0.12)
 & vidutinis, 4\newline{\footnotesize DET 1–4\newline FTR 3+} \\
\midrule
P3.4 & Rezervacijos atšaukimas\newline EI
 & \textbf{3}: pasirinkta rezervacija, mygtukas „Atšaukti“, pranešimas
 & \textbf{FTR 2}: Rezervacijos ir lankomumas (būsena „atšaukta“, atšaukimo
   laikas, kas atšaukė); Užsiėmimas (ar dar neprasidėjo)
 & žemas, 3\newline{\footnotesize DET 1–4\newline FTR 2} \\
\midrule
P3.5 & Dalyvio priskyrimas užsiėmimui (treneris, administratorius)\newline EI
 & \textbf{4}: užsiėmimas, dalyvis (įvedamas jo el.~pašto adresas; sistema
   randa paskyrą), mygtukas „Priskirti“, pranešimas
 & \textbf{FTR 3}: Užsiėmimas (vietų skaičius, ar jį veda šis treneris);
   Rezervacijos ir lankomumas (užimtos vietos, nauja rezervacija, kas
   užregistravo); Naudotojo paskyra (dalyvio paskyra pagal el.~paštą, ar ji
   aktyvi)
 & vidutinis, 4\newline{\footnotesize DET 1–4\newline FTR 3+} \\
\midrule
P3.6 & Dalyvio pašalinimas iš užsiėmimo\newline EI
 & \textbf{4}: užsiėmimas, pasirinktas dalyvis, mygtukas „Pašalinti“,
   pranešimas
 & \textbf{FTR 2}: Rezervacijos ir lankomumas (būsena „pašalinta“,
   pašalinimo laikas, kas pašalino); Užsiėmimas (ar jį veda šis treneris)
 & žemas, 3\newline{\footnotesize DET 1–4\newline FTR 2} \\
\midrule
P3.7 & Faktinio lankomumo pažymėjimas ir taisymas\newline EI
 & \textbf{5}: užsiėmimas, dalyvis, lankomumo žyma (dalyvavo arba
   nedalyvavo), mygtukas „Išsaugoti“, pranešimas
 & \textbf{FTR 2}: Rezervacijos ir lankomumas (užsiregistravę dalyviai,
   lankomumo įrašas); Užsiėmimas (ar jau prasidėjo, ar jį veda šis treneris)
 & vidutinis, 4\newline{\footnotesize DET 5–15\newline FTR 2} \\
\midrule
P3.8 & Dalyvio rezervacijų peržiūra\newline EQ
 & \textbf{7}: filtras „būsimos arba praėjusios“, užsiėmimo pavadinimas,
   data, pradžios laikas, vieta, treneris, rezervacijos būsena
 & \textbf{FTR 4}: Rezervacijos ir lankomumas (būsena); Užsiėmimas
   (pavadinimas, data, laikas); Sporto erdvė (vieta); Naudotojo paskyra
   (treneris)
 & aukštas, 6\newline{\footnotesize DET 5–15\newline FTR 3+} \\
\midrule
P3.9 & Užsiregistravusių dalyvių sąrašas\newline EQ
 & \textbf{6}: užsiėmimas (pasirinktas; antraštėje rodomas jo pavadinimas),
   dalyvio vardas, pavardė (atskiri stulpeliai), telefonas, rezervacijos
   būsena, rezervavimo laikas
 & \textbf{FTR 3}: Užsiėmimas (pasirinkto užsiėmimo pavadinimas; ar jį veda
   šis treneris); Rezervacijos ir lankomumas (būsena, rezervavimo laikas);
   Naudotojo paskyra (vardas, pavardė, telefonas)
 & aukštas, 6\newline{\footnotesize DET 5–15\newline FTR 3+} \\
\midrule
P3.10 & Dalyvio lankomumo istorija (dalyviui ir treneriui)\newline EQ
 & \textbf{8}: dalyvis (vardas ir pavardė antraštėje; treneris jį
   pasirenka), laikotarpis, užsiėmimo pavadinimas, režimas, data, pradžios
   laikas, rezervacijos būsena, lankomumo žyma
 & \textbf{FTR 3}: Rezervacijos ir lankomumas (būsena, lankomumo žyma);
   Užsiėmimas (pavadinimas, režimas, data, laikas); Naudotojo paskyra
   (dalyvio vardas ir pavardė)
 & aukštas, 6\newline{\footnotesize DET 5–15\newline FTR 3+} \\
\midrule
P3.11 & Priminimas apie artėjantį užsiėmimą\newline EO
 & \textbf{6}: dalyvio el.~pašto adresas, vardas, data, laikas, užsiėmimo
   pavadinimas, vieta. Laiško tekstas: „[Vardas], primename: rytoj [data]
   [laikas]~– [užsiėmimas], [vieta]“
 & \textbf{FTR 4}: Rezervacijos ir lankomumas (rytdienos aktyvios
   rezervacijos, žyma „priminimas išsiųstas“); Užsiėmimas (pavadinimas,
   data, laikas); Sporto erdvė (vieta); Naudotojo paskyra (vardas,
   el.~paštas)
 & aukštas, 7\newline{\footnotesize DET 6–19\newline FTR 4+} \\
\midrule
P3.12 & Trenerio dalyvių sąrašas ir paieška\newline EQ
 & \textbf{6}: dalyvio vardas, pavardė (pagal juos ir ieškoma), telefonas,
   el.~paštas, plaukimo lygis, mygtukas „Ieškoti“
 & \textbf{FTR 3}: Užsiėmimas (trenerio užsiėmimai); Rezervacijos ir
   lankomumas (kas į juos užsiregistravo); Naudotojo paskyra (vardas,
   pavardė, telefonas, el.~paštas, plaukimo lygis)
 & aukštas, 6\newline{\footnotesize DET 5–15\newline FTR 3+} \\
\midrule
P3.13 & Dalyvio kortelės peržiūra (treneriui)\newline EQ
 & \textbf{10}: dalyvio pasirinkimas sąraše P3.12 (veiksmas, kuriuo
   atidaroma kortelė), vardas, pavardė, gimimo data, telefonas, el.~paštas,
   plaukimo lygis, tikslai, pastabos treneriui, paskyros sukūrimo data
 & \textbf{FTR 1}: Naudotojo paskyra (naudotojo duomenys ir dalyvio
   profilis)
 & žemas, 3\newline{\footnotesize DET 5–15\newline FTR 0–1} \\
\end{longtable}
\endgroup

\paragraph{Kodėl taip skaičiuota.}
Toliau paaiškinti sprendimai, kurie lėmė \ref{tab:fp-p3}~lentelės skaičius.

\begin{itemize}
  \item Failas „Rezervacijos ir lankomumas“ yra atskiras ILF, o ne
    užsiėmimo RET. Rezervacija sieja du nepriklausomus failus~– dalyvio
    paskyrą ir užsiėmimą~– ir nepriklauso nė vienam iš jų; naudotojas ją
    mato kaip atskirą objektą („Mano rezervacijos“, lankomumo istorija per
    visus užsiėmimus), o įrašai išlieka ir atšaukus užsiėmimą (P2.5) ar
    deaktyvavus paskyrą (P0.12), nes iš jų sudaroma istorija (P3.10) ir
    ataskaita (P7.2). Paskaitos parkavimo pavyzdyje būsenos failas,
    siejantis automobilį ir zoną, taip pat yra atskiras ILF. Kliento
    įvertinimas, priešingai, priklauso vienai asmeninei programai, todėl
    jis yra RET (P4.1). Faile dar saugoma, kas užregistravo, kas ir kada
    atšaukė ar pašalino rezervaciją, ar išsiųstas priminimas, kas ir kada
    pažymėjo lankomumą ir kada jį pataisė; \ref{tab:fp-p3}~lentelėje nurodyta,
    kuri funkcija šiuos laukus įrašo. Net be šių laukų failas liktų žemo
    sudėtingumo (DET~1–19).
  \item Tvarkaraštis (P3.2) yra EO, o ne EQ, nes laisvų vietų skaičius
    nesaugomas: jis kaskart apskaičiuojamas iš užsiėmimo vietų skaičiaus
    atimant aktyvių rezervacijų skaičių. Išvestis su skaičiavimais yra EO,
    o EQ tik pateikia saugomus duomenis. Užsiėmimo pavadinimas įvedamas
    planuojant užsiėmimą (P2.2, P2.3) ir saugomas faile „Užsiėmimas“.
  \item Rezervavimas (P3.3) ir priskyrimas (P3.5), taip pat atšaukimas
    (P3.4) ir pašalinimas (P3.6) yra atskiros funkcijos, nors keičia tą
    patį failą. Priskirdamas ar šalindamas treneris arba administratorius
    nurodo dalyvį (papildomas DET), o kai veiksmą atlieka treneris, sistema
    dar tikrina, ar užsiėmimą veda jis; dalyvis rezervuoja ir atšaukia savo
    vardu, todėl jo tapatybė DET nelaikoma.
  \item Lankomumo pažymėjimas ir taisymas (P3.7) yra viena įvestis: tas
    pats lankomumo ekranas, tie patys laukai, failai ir patikros; taisant
    tik perrašoma žyma ir įrašomas pataisymo laikas.
  \item Lankomumo istorija (P3.10) dalyviui ir treneriui yra viena
    užklausa: abiem atvejais rodomi tie patys laukai (įskaitant dalyvio
    vardą ir pavardę antraštėje), skaitomi tie patys trys failai ir
    taikoma ta pati logika. Skiriasi tik tai, kaip parenkamas dalyvis:
    treneris jį pasirenka, o dalyviui rodoma jo paties istorija. Todėl
    skaičiuojama 8~DET ir 3~FTR.
  \item Panašios funkcijos skaičiuojamos atskirai, nes skiriasi jų DET ir
    FTR. Rezervacijų peržiūra (P3.8) rodo vietą ir trenerį, todėl skaito ir
    failą „Sporto erdvė“, o iš failo „Naudotojo paskyra“~– trenerio vardą
    (4~FTR);
    lankomumo istorija (P3.10) vietos ir trenerio nerodo, bet rodo
    užsiėmimo režimą ir lankomumo žymą (3~FTR). Trenerio dalyvių sąrašas
    (P3.12) skiriasi nuo administratoriaus naudotojų sąrašo (P0.13), kuris
    skaito tik paskyras (1~FTR) ir rodo visus naudotojus, ir nuo trenerio
    klientų sąrašo (P4.9), kuris rodo tik klientus, turinčius asmeninę
    programą, ir skaito failą „Asmeninė programa“.
  \item Priminimas (P3.11) yra EO, o ne EQ: sistema ne tik pateikia
    duomenis, bet ir taiko verslo logiką bei keičia vidinį failą~– kasdien
    pati atrenka rytdienos aktyvias rezervacijas, išsiunčia laišką dalyvio
    el.~paštu ir pažymi „priminimas išsiųstas“, kad laiškas nebūtų
    išsiųstas dukart. EQ tik pateikia informaciją ir vidinių failų
    nekeičia. Tai ir ne EI, nes per sistemos ribą duomenys neįeina, o
    pagrindinis veiksmo tikslas~– pranešti dalyviui.
  \item Trenerio dalyviai yra jo užsiėmimų dalyviai: dalyvis pasirenka
    trenerį, kai rezervuoja jo užsiėmimą (P3.3), o individualiems
    užsiėmimams dalyvį trenerio užsiėmimui priskiria administratorius arba
    pats treneris (P3.5). Todėl atskiros dalyvio priskyrimo treneriui
    funkcijos nėra, o sąrašas P3.12 dalyvius randa per trenerio užsiėmimus
    ir jų rezervacijas (3~FTR).
  \item Dalyvio kortelė (P3.13) ir naudotojo paskyros peržiūra (P0.14)
    skaito tą patį failą, bet rodo skirtingus laukus: treneris mato sportui
    reikalingus profilio duomenis (plaukimo lygį, tikslus, pastabas
    treneriui), o administratorius~– paskyros duomenis (prisijungimo
    duomenis, vaidmenį, būseną). Šiuos profilio laukus dalyvis pildo
    savo profilyje (P0.8). Kortelė
    atidaroma tik iš trenerio sąrašo P3.12, kuriame jau yra tik trenerio
    dalyviai, todėl ji, kaip ir P0.14, skaito tik paskyrą (1~FTR). DET
    rinkiniai skiriasi, todėl tai dvi atskiros užklausos.
\end{itemize}
